% GENERATED FILE. Edit canonical lessons, then run npm run generate:books.
% canonical-source-hash: fnv1a64:4ade86015c18b65c
% canonical-lessons: TE-C96-surname, TE-C96-nickname, TE-C96-age, TE-C96-birthday, TE-C96-signature

\chapter{Names, Ages, Birthdays}
\label{ch:te-names-ages-birthdays}
\hlchaptermodality{96}

\begin{chapteropening}
\textbf{By the end of this chapter:} \emph{I can give a surname and a nickname, ask an age, talk about a birthday, and sign my name.}

Complete TE-C96-signature: the details of a name, from surname to signature.
\end{chapteropening}

\section[\texorpdfstring{\te{ఇంటిపేరు} (iṇṭipēru)}{ఇంటిపేరు (iṇṭipēru)}]{\te{ఇంటిపేరు} (iṇṭipēru) — a surname, the family name}

\label{lesson:TE-C96-surname}



\begin{quote}
Before the new one: say the Telugu for please show me, then the Telugu for I could not hear it.
\end{quote}

\subsection*{You'll want to know: \te{ఇంటిపేరు}}
\textbf{\te{ఇంటిపేరు}} (\emph{iṇṭipēru}) — \textquotedblleft{}a surname\textquotedblright{}, literally \textquotedblleft{}the house name\textquotedblright{}.

It is \textbf{\te{ఇంటి}}, \textquotedblleft{}of the house\textquotedblright{}, and \textbf{\te{పేరు}}, \textquotedblleft{}name\textquotedblright{}. In Telugu the house name comes FIRST, before the given name, so a form that asks for a first name is asking the wrong question.

\begin{cousinweb}
\textbf{\te{ఇంటిపేరు}} (\emph{iṇṭipēru}) is native Dravidian: \emph{iṇṭi}, \textquotedblleft{}of the house\textquotedblright{} (DEDR 494), and \emph{pēru}, \textquotedblleft{}name\textquotedblright{} (DEDR 4410).
\end{cousinweb}

\begin{grammarlens}[title={What you've built}]
The family name, and where it goes.
\end{grammarlens}

\subsection*{Your turn}
\begin{itemize}
\raggedright
  \item \emph{Say it:} \emph{iṇṭipēru}
  \item \emph{Say it:} \emph{iṇṭipēru}, once more
  \item \emph{Say it:} say \emph{nā iṇṭipēru ...}, then your given name
  \item \emph{Recall:} say the Telugu for please show me, then the Telugu for I could not hear it, then say \emph{iṇṭipēru} again
\end{itemize}

\begin{tcolorbox}[breakable,colback=teal!4,colframe=teal!35!black,title={Before you move on}]
What does \te{ఇంటిపేరు} mean? (\textquotedblleft{}A surname, the family name\textquotedblright{}.) Say it once more.
\end{tcolorbox}
\section[\texorpdfstring{\te{ముద్దుపేరు} (muddupēru)}{ముద్దుపేరు (muddupēru)}]{\te{ముద్దుపేరు} (muddupēru) — a nickname, a pet name}

\label{lesson:TE-C96-nickname}



\begin{quote}
Before the new one: say the Telugu for I could not hear it, then the Telugu for a surname.
\end{quote}

\subsection*{You'll want to know: \te{ముద్దుపేరు}}
\textbf{\te{ముద్దుపేరు}} (\emph{muddupēru}) — \textquotedblleft{}a nickname, a pet name\textquotedblright{}.

\textbf{\te{ముద్దు}} is \textquotedblleft{}a kiss, affection\textquotedblright{}, so a nickname is an \emph{affection name}: the one family uses.

\begin{cousinweb}
\textbf{\te{ముద్దుపేరు}} (\emph{muddupēru}) is native Dravidian: \emph{muddu}, \textquotedblleft{}kiss\textquotedblright{}, one word with Tamil \emph{muttu} (DEDR 4960), and \emph{pēru}, \textquotedblleft{}name\textquotedblright{} (DEDR 4410). DEDR compares Prakrit \emph{muddī}, \textquotedblleft{}kiss\textquotedblright{}.
\end{cousinweb}

\begin{grammarlens}[title={What you've built}]
The name the family uses.
\end{grammarlens}

\subsection*{Your turn}
\begin{itemize}
\raggedright
  \item \emph{Say it:} \emph{muddupēru}
  \item \emph{Say it:} \emph{muddupēru}, once more
  \item \emph{Say it:} say \emph{nā muddupēru ...}
  \item \emph{Recall:} say the Telugu for I could not hear it, then the Telugu for a surname, then say \emph{muddupēru} again
\end{itemize}

\begin{tcolorbox}[breakable,colback=teal!4,colframe=teal!35!black,title={Before you move on}]
What does \te{ముద్దుపేరు} mean? (\textquotedblleft{}A nickname, a pet name\textquotedblright{}.) Say it once more.
\end{tcolorbox}
\section[\texorpdfstring{\te{వయసు} (vayasu)}{వయసు (vayasu)}]{\te{వయసు} (vayasu) — age}

\label{lesson:TE-C96-age}



\begin{quote}
Before the new one: say the Telugu for a surname, then the Telugu for a nickname.
\end{quote}

\subsection*{You'll want to know: \te{వయసు}}
\textbf{\te{వయసు}} (\emph{vayasu}) — \textquotedblleft{}age\textquotedblright{}.

You have asked it already, inside \textbf{\te{నీ} \te{వయసెంత}?}. Now the word stands alone: \textbf{\te{మీ} \te{వయసు} \te{ఎంత}?} — \textquotedblleft{}what is your age?\textquotedblright{}, the respectful form.

\begin{cousinweb}
\textbf{\te{వయసు}} (\emph{vayasu}) is Sanskrit \emph{vayas}, \textquotedblleft{}prime of life, age\textquotedblright{}, taken over whole as a learned loan, a tatsama, with Telugu's added \emph{-u} (CDIAL 11305).
\end{cousinweb}

\begin{grammarlens}[title={What you've built}]
The word you have been using inside a question, taken out.
\end{grammarlens}

\subsection*{Your turn}
\begin{itemize}
\raggedright
  \item \emph{Say it:} \emph{vayasu}
  \item \emph{Say it:} \emph{vayasu}, once more
  \item \emph{Say it:} ask \emph{mī vayasu enta?}
  \item \emph{Recall:} say the Telugu for a surname, then the Telugu for a nickname, then say \emph{vayasu} again
\end{itemize}

\begin{tcolorbox}[breakable,colback=teal!4,colframe=teal!35!black,title={Before you move on}]
What does \te{వయసు} mean? (\textquotedblleft{}Age\textquotedblright{}.) Say it once more.
\end{tcolorbox}
\section[\texorpdfstring{\te{పుట్టినరోజు} (puṭṭinarōju)}{పుట్టినరోజు (puṭṭinarōju)}]{\te{పుట్టినరోజు} (puṭṭinarōju) — a birthday}

\label{lesson:TE-C96-birthday}



\begin{quote}
Before the new one: say the Telugu for a nickname, then the Telugu for age.
\end{quote}

\subsection*{You'll want to know: \te{పుట్టినరోజు}}
\textbf{\te{పుట్టినరోజు}} (\emph{puṭṭinarōju}) — \textquotedblleft{}a birthday\textquotedblright{}, literally \textquotedblleft{}the day one was born\textquotedblright{}.

\textbf{\te{రోజు}} is the day you know. On the day, say \textbf{\te{పుట్టినరోజు} \te{శుభాకాంక్షలు}} — \textquotedblleft{}happy birthday\textquotedblright{}, with the good wishes you already have.

\begin{cousinweb}
\textbf{\te{పుట్టినరోజు}} (\emph{puṭṭinarōju}) is native \emph{puṭṭina}, \textquotedblleft{}on which one was born\textquotedblright{}, from \emph{puṭṭu}, \textquotedblleft{}to be born\textquotedblright{} (DEDR 4264), and \emph{rōju}, \textquotedblleft{}day\textquotedblright{}, from Persian \emph{rōz}.
\end{cousinweb}

\begin{grammarlens}[title={What you've built}]
A day to name, and a wish to make on it.
\end{grammarlens}

\subsection*{Your turn}
\begin{itemize}
\raggedright
  \item \emph{Say it:} \emph{puṭṭinarōju}
  \item \emph{Say it:} \emph{puṭṭinarōju}, once more
  \item \emph{Say it:} say \emph{puṭṭinarōju śubhākāṅkṣalu}
  \item \emph{Recall:} say the Telugu for a nickname, then the Telugu for age, then say \emph{puṭṭinarōju} again
\end{itemize}

\begin{tcolorbox}[breakable,colback=teal!4,colframe=teal!35!black,title={Before you move on}]
What does \te{పుట్టినరోజు} mean? (\textquotedblleft{}A birthday\textquotedblright{}.) Say it once more.
\end{tcolorbox}
\section[\texorpdfstring{\te{సంతకం} (santakaṁ)}{సంతకం (santakaṁ)}]{\te{సంతకం} (santakaṁ) — a signature}

\label{lesson:TE-C96-signature}



\begin{quote}
Before the new one: say the Telugu for age, then the Telugu for a birthday.
\end{quote}

\subsection*{You'll want to know: \te{సంతకం}}
\textbf{\te{సంతకం}} (\emph{santakaṁ}) — \textquotedblleft{}a signature\textquotedblright{}.

\textbf{\te{ఇక్కడ} \te{సంతకం} \te{చేయండి}} — \textquotedblleft{}please sign here\textquotedblright{}. A signature is something you \emph{do}, with the \textbf{\te{చేయు}} of \textbf{\te{పని} \te{చేయు}}.

\begin{cousinweb}
The origin of \textbf{\te{సంతకం}} (\emph{santakaṁ}) is uncertain. Neither DEDR nor CDIAL records it.
\end{cousinweb}

\begin{grammarlens}[title={What you've built}]
Surname, nickname, age, birthday, signature: every box on a form.
\end{grammarlens}

\subsection*{Your turn}
\begin{itemize}
\raggedright
  \item \emph{Say it:} \emph{santakaṁ}
  \item \emph{Say it:} \emph{santakaṁ}, once more
  \item \emph{Say it:} say \emph{ikkaḍa santakaṁ cēyaṇḍi}
  \item \emph{Recall:} say the Telugu for age, then the Telugu for a birthday, then say \emph{santakaṁ} again
\end{itemize}

\begin{tcolorbox}[breakable,colback=teal!4,colframe=teal!35!black,title={Before you move on}]
What does \te{సంతకం} mean? (\textquotedblleft{}A signature\textquotedblright{}.) Say it once more.
\end{tcolorbox}
